% HQIV Topology / TUFT Hopf bridge: finite approximations, contact spectra,
% and the discrete--continuous interface.
% pdflatex
\documentclass[11pt]{article}

\usepackage{amsmath,amssymb,amsthm,mathtools}
\usepackage{enumitem}
\usepackage[margin=1in]{geometry}
\usepackage{hyperref}
\hypersetup{colorlinks=true,linkcolor=blue,citecolor=blue,urlcolor=blue}
\usepackage[final]{microtype}
\setlength{\emergencystretch}{3em}

\newcommand{\leanrepo}{https://github.com/HQIV/hqiv-lean/blob/main}
% Lean module / identifier helpers: visible text is rendered with \nolinkurl
% so that slashes, underscores, and dots become legal break points.
\newcommand{\leanfile}[1]{\href{\leanrepo/#1}{\nolinkurl{#1}}}
\newcommand{\leanid}[1]{\nolinkurl{#1}}

\newtheorem{theorem}{Theorem}[section]
\newtheorem{lemma}[theorem]{Lemma}
\newtheorem{proposition}[theorem]{Proposition}
\newtheorem{corollary}[theorem]{Corollary}
\newtheorem{definition}[theorem]{Definition}
\newtheorem{remark}[theorem]{Remark}

\title{Topology, the Complex Hopf Fibration, and Finite Approximations in HQIV:\\[0.35em]
TUFT Alignments, Contact Spectra, and the Discrete--Continuous Interface}

\author{Steven Ettinger Jr\\
{\small Independent Researcher --- \texttt{steven@disregardfiat.tech}}}
\date{Preprint v1 --- \today}

\begin{document}
\maketitle

\begin{abstract}
\noindent\textbf{Position in the publication chain.}
This is a Tier-2 topology companion in the HQIV preprint series. It supplies
the topological ``why'' layer for structures already present in the HQIV
discrete null-lattice and Fano/octonion carrier: the rigidity of three
generations, the shell-surface readouts, the effective continuous $\xi$
chart, and the informational-energy mass row. The upstream anchors are the
main HQIV framework \cite{ettinger2026hqiv}, the light-cone null-lattice
readout \cite{hqiv-lightcone-paper}, and the $\mathfrak{so}(8)$ closure
\cite{hqiv-so8-paper}. The Standard-Model Lagrangian paper
the combined TUFT+SM synthesis \cite{hqiv-tuft-sm-lagrangian-paper} is the intended Tier-2 foundation
that will cross-cite this work once all records carry minted DOIs. The
three published Tier-1 companions \cite{hqiv-3d-growth-paper,hqiv-oct-action-paper,hqiv-kirchhoff-paper}
provide the immediate dimensional, variational, and finite-mode spine.

\noindent\textbf{Result.}
Jenny Lorraine Nielsen's Topological Unified Field Theory (TUFT) on the
complex Hopf fibration \cite{NielsenTUFT2026} proves that any U(1)-complete
gauge theory with discrete charge quantization must, up to homotopy and
bundle isomorphism, be formulated on the universal Hopf fibration
$S^1\to S^\infty\to\mathbb{CP}^\infty$ and its finite approximations
$S^1\to S^{2n+1}\to\mathbb{CP}^n$. The Standard-Model gauge factors and a
gravity sector with fiber-induced torsion arise as reductions along nested
Hopf shells, with a contact Beltrami spectrum on each shell and
zeta-regularized sector determinants supplying intrinsic mass scales.
This paper records the precise, machine-checked alignments already present
in \texttt{HQIV\_LEAN} between TUFT's finite-shell contact geometry and
HQIV's discrete null-lattice + Fano carrier, together with the explicit
chart distinctions that keep the two frameworks from being silently
identified. The 4/3 lock-in-neighbour match, the three-generation
count from integrable fibre windings, the Beltrami versus Peter--Weyl
normalisation distinction, and the placement of the 3/2 ratio in Fano
holonomy rows rather than geometric resonance steps have all been
formally verified in Lean without the use of `sorry`. These alignments are bidirectional: TUFT supplies a
topological derivation for structures HQIV had taken as geometric axioms;
HQIV's patch ontology supplies the finite-truncation discipline that
prevents either framework from treating $\mathbb{CP}^\infty$ or the
continuous $\xi$ chart as literal ontological commitments.

\noindent\textbf{Patch-ontology reader contract.}
Per the patch-theory reader contract imported below, the discrete patch
layer (shell indices $m\in\mathbb{N}$, finite index types $\mathrm{Fin}\,n$,
horizon-limited charts) is the ontology. Smooth-manifold notation,
continuous $\xi$ charts, and $\mathbb{CP}^n$ language appearing in this
paper are translation layers for comparison with the TUFT literature and
with standard physics. They are not claims that a fundamental continuum
spacetime or an ultraviolet limit has been recovered or is required.
Taking a literal $m\to\infty$ or sub-Planck continuum limit on the load-bearing
objects (lock-in at $\mathrm{referenceM}=4$, Fano holonomy rows, resonance
steps, informational-energy localization) would erase the discrete index
that does the predictive work.

\noindent\textbf{Status.}
The alignments recorded here, including the typed
\hyperref[lean:hopf-shell-complex]{Hopf-shell complex + ContactBeltrami}
layer (now part of the official `HQIVParallelPoincare` and `HQIVPhysics`
build globs and verified green), the three-row phase-contribution assembler with
admissible-cycle and orientation fields, the canonical matrix-torsion (phase-lift)
construction with skew-adjoint action, the concrete three-shell non-factorability witness, and
especially the exact algebraic trapped-Casimir factorization of the
binding cell together with the active per-shell curvature-imprint program,
have been formally verified or constructed in Lean without the use of
`sorry`. A dedicated section below articulates the principal contributions of
Nielsen's framework (the classifying-space necessity of the complex Hopf
fibration for charge-quantized theories, the finite-approximation
discipline, the contact Beltrami spectrum and fibre holonomy phases as a
potential source for mixing phenomena, and the single-scale derivation
programme, including a re-interpretation of apparent dark-sector effects
as internal to HQIV's dimensioned quantities at the present epoch). The
paper further records the precise points at which the patch ontology of
HQIV and the Lean formalization supply mutual reinforcement. Certain
deeper alignments remain the subject of ongoing formalization efforts and
are catalogued in the accompanying technical note on formalization
targets. No universality theorem or single-Fermi-constant derivation is
imported from Nielsen's work.
Reader-friendly anchors live in Appendix~\ref{app:lean-index}.
\end{abstract}

% Shared reader contract: patch QFT + discrete numerics.
% From papers/*.tex:     % Shared reader contract: patch QFT + discrete numerics.
% From papers/*.tex:     % Shared reader contract: patch QFT + discrete numerics.
% From papers/*.tex:     \input{include/patch_theory_messaging.tex}
% From papers/paper/*.tex: \input{../include/patch_theory_messaging.tex}

\paragraph{Patch QFT and discrete numerics (reader contract).}
HQIV (Horizon-Quantized Informational Vacuum) is formulated as a \emph{patch theory}: quantum fields, actions, and physical readouts are attached to discrete patches---shell indices $m\in\mathbb{N}$, finite spacetime index types (e.g.\ $\mathrm{Fin}\,4$), and horizon-limited charts---rather than to a hypothetical fundamental smooth spacetime obtained as a mesh-refinement or ultraviolet limit.
Accordingly, when this text refers to ``QFT,'' it means \emph{patch QFT}: patching rules, finite measurement ledgers, and variational identities on the discrete spine; it is not the claim that the theory is \emph{defined} by recovering ordinary continuum quantum field theory through such a limit.
The numbers that admit the sharpest statements---mass and coupling ladders, curvature ratios, shell thermodynamics, and rational witnesses in \texttt{HQIV\_LEAN}---are objects of \emph{discrete mathematics} on $\mathbb{N}$ (combinatorial growth, cumulative simplex ledgers) together with the ordered field $\mathbb{R}$ as used in the proof assistant for analysis, bounds, and phenomenological display.
Smooth-manifold or continuum field notation, when it appears, is a \emph{translation layer} for comparison with standard physics literature; it does not supersede the discrete definitions.

\paragraph{A continuum is not required, and in places would destroy these results.}
The discrete patch layer is \emph{not} a numerical approximation to a more fundamental continuum theory awaiting future construction; it is the ontology.  Where the literature treats ``the continuum limit'' as the goal, HQIV treats it as a \emph{calculation approximation} for comparison with standard formulas.  The continuum form is an optional convenience for comparison; the proved discrete statement is what carries the physics.  In several load-bearing places---the curvature imprint $\delta_E(m)$ at shell $m$, the harmonic ladder masses $m_\tau\to m_\mu\to m_e$, the lock-in vev $v=\sqrt{\eta_{\rm paper}\,\Omega_k(m_{\rm lockin};m_{\rm lockin})}$, the proton anchor at $m=\mathrm{referenceM}=4$, the baryon-asymmetry $\eta\propto 6^7\sqrt{3}/(m+1)$, and the rational coefficients $\alpha=3/5$, $\gamma=2/5$, $\alpha_{\rm GUT}=1/42$---taking a literal $m\to\infty$ or sub-Planck continuum limit would \emph{erase} the discrete index that is doing the predictive work.  Continuum forms of these objects are therefore not preferred refinements; they are degenerate, information-losing readouts of the same patch data.

\paragraph{An exception: $\mathfrak{so}(8)$ as a de-facto continuum algebra from a discrete construction.}
Principled exception to the discrete-only stance: the closed Lie algebra $\mathfrak{so}(8) \supseteq G_2\cup\{\Delta\}$ generated from the Fano octonion left-multiplication table and the phase-lift generator $\Delta\in(\mathrm{e}_1,\mathrm{e}_7)$.  Although $\mathfrak{so}(8)$ is a continuous (28-dimensional real) Lie algebra, it is built here entirely from \emph{discrete} ingredients (the seven imaginary octonion units, a finite Fano table, and one $\pi/2$ rotation), and its closure to 28 dimensions is a finite linear-algebra certificate (see the \texttt{Hqiv.SO8Closure} / \texttt{Hqiv.SO8ClosureInterface} stack).  The Lie-group structure is therefore a \emph{de-facto continuum algebra obtained from a discrete construction}: continuous in parameter space, but with no continuum spacetime, no continuum measure, and no UV limit attached.  When this manuscript or its companions write $\exp(t\,\Delta)\in\mathrm{SO}(8)$ or treat $\mathfrak{so}(8)$ rotations as smooth, those are statements internal to a finite-dimensional Lie group whose generators are certified by discrete data; they are \emph{not} a continuum-spacetime commitment and do not require a continuum-QFT scaffolding.

% From papers/paper/*.tex: % Shared reader contract: patch QFT + discrete numerics.
% From papers/*.tex:     \input{include/patch_theory_messaging.tex}
% From papers/paper/*.tex: \input{../include/patch_theory_messaging.tex}

\paragraph{Patch QFT and discrete numerics (reader contract).}
HQIV (Horizon-Quantized Informational Vacuum) is formulated as a \emph{patch theory}: quantum fields, actions, and physical readouts are attached to discrete patches---shell indices $m\in\mathbb{N}$, finite spacetime index types (e.g.\ $\mathrm{Fin}\,4$), and horizon-limited charts---rather than to a hypothetical fundamental smooth spacetime obtained as a mesh-refinement or ultraviolet limit.
Accordingly, when this text refers to ``QFT,'' it means \emph{patch QFT}: patching rules, finite measurement ledgers, and variational identities on the discrete spine; it is not the claim that the theory is \emph{defined} by recovering ordinary continuum quantum field theory through such a limit.
The numbers that admit the sharpest statements---mass and coupling ladders, curvature ratios, shell thermodynamics, and rational witnesses in \texttt{HQIV\_LEAN}---are objects of \emph{discrete mathematics} on $\mathbb{N}$ (combinatorial growth, cumulative simplex ledgers) together with the ordered field $\mathbb{R}$ as used in the proof assistant for analysis, bounds, and phenomenological display.
Smooth-manifold or continuum field notation, when it appears, is a \emph{translation layer} for comparison with standard physics literature; it does not supersede the discrete definitions.

\paragraph{A continuum is not required, and in places would destroy these results.}
The discrete patch layer is \emph{not} a numerical approximation to a more fundamental continuum theory awaiting future construction; it is the ontology.  Where the literature treats ``the continuum limit'' as the goal, HQIV treats it as a \emph{calculation approximation} for comparison with standard formulas.  The continuum form is an optional convenience for comparison; the proved discrete statement is what carries the physics.  In several load-bearing places---the curvature imprint $\delta_E(m)$ at shell $m$, the harmonic ladder masses $m_\tau\to m_\mu\to m_e$, the lock-in vev $v=\sqrt{\eta_{\rm paper}\,\Omega_k(m_{\rm lockin};m_{\rm lockin})}$, the proton anchor at $m=\mathrm{referenceM}=4$, the baryon-asymmetry $\eta\propto 6^7\sqrt{3}/(m+1)$, and the rational coefficients $\alpha=3/5$, $\gamma=2/5$, $\alpha_{\rm GUT}=1/42$---taking a literal $m\to\infty$ or sub-Planck continuum limit would \emph{erase} the discrete index that is doing the predictive work.  Continuum forms of these objects are therefore not preferred refinements; they are degenerate, information-losing readouts of the same patch data.

\paragraph{An exception: $\mathfrak{so}(8)$ as a de-facto continuum algebra from a discrete construction.}
Principled exception to the discrete-only stance: the closed Lie algebra $\mathfrak{so}(8) \supseteq G_2\cup\{\Delta\}$ generated from the Fano octonion left-multiplication table and the phase-lift generator $\Delta\in(\mathrm{e}_1,\mathrm{e}_7)$.  Although $\mathfrak{so}(8)$ is a continuous (28-dimensional real) Lie algebra, it is built here entirely from \emph{discrete} ingredients (the seven imaginary octonion units, a finite Fano table, and one $\pi/2$ rotation), and its closure to 28 dimensions is a finite linear-algebra certificate (see the \texttt{Hqiv.SO8Closure} / \texttt{Hqiv.SO8ClosureInterface} stack).  The Lie-group structure is therefore a \emph{de-facto continuum algebra obtained from a discrete construction}: continuous in parameter space, but with no continuum spacetime, no continuum measure, and no UV limit attached.  When this manuscript or its companions write $\exp(t\,\Delta)\in\mathrm{SO}(8)$ or treat $\mathfrak{so}(8)$ rotations as smooth, those are statements internal to a finite-dimensional Lie group whose generators are certified by discrete data; they are \emph{not} a continuum-spacetime commitment and do not require a continuum-QFT scaffolding.


\paragraph{Patch QFT and discrete numerics (reader contract).}
HQIV (Horizon-Quantized Informational Vacuum) is formulated as a \emph{patch theory}: quantum fields, actions, and physical readouts are attached to discrete patches---shell indices $m\in\mathbb{N}$, finite spacetime index types (e.g.\ $\mathrm{Fin}\,4$), and horizon-limited charts---rather than to a hypothetical fundamental smooth spacetime obtained as a mesh-refinement or ultraviolet limit.
Accordingly, when this text refers to ``QFT,'' it means \emph{patch QFT}: patching rules, finite measurement ledgers, and variational identities on the discrete spine; it is not the claim that the theory is \emph{defined} by recovering ordinary continuum quantum field theory through such a limit.
The numbers that admit the sharpest statements---mass and coupling ladders, curvature ratios, shell thermodynamics, and rational witnesses in \texttt{HQIV\_LEAN}---are objects of \emph{discrete mathematics} on $\mathbb{N}$ (combinatorial growth, cumulative simplex ledgers) together with the ordered field $\mathbb{R}$ as used in the proof assistant for analysis, bounds, and phenomenological display.
Smooth-manifold or continuum field notation, when it appears, is a \emph{translation layer} for comparison with standard physics literature; it does not supersede the discrete definitions.

\paragraph{A continuum is not required, and in places would destroy these results.}
The discrete patch layer is \emph{not} a numerical approximation to a more fundamental continuum theory awaiting future construction; it is the ontology.  Where the literature treats ``the continuum limit'' as the goal, HQIV treats it as a \emph{calculation approximation} for comparison with standard formulas.  The continuum form is an optional convenience for comparison; the proved discrete statement is what carries the physics.  In several load-bearing places---the curvature imprint $\delta_E(m)$ at shell $m$, the harmonic ladder masses $m_\tau\to m_\mu\to m_e$, the lock-in vev $v=\sqrt{\eta_{\rm paper}\,\Omega_k(m_{\rm lockin};m_{\rm lockin})}$, the proton anchor at $m=\mathrm{referenceM}=4$, the baryon-asymmetry $\eta\propto 6^7\sqrt{3}/(m+1)$, and the rational coefficients $\alpha=3/5$, $\gamma=2/5$, $\alpha_{\rm GUT}=1/42$---taking a literal $m\to\infty$ or sub-Planck continuum limit would \emph{erase} the discrete index that is doing the predictive work.  Continuum forms of these objects are therefore not preferred refinements; they are degenerate, information-losing readouts of the same patch data.

\paragraph{An exception: $\mathfrak{so}(8)$ as a de-facto continuum algebra from a discrete construction.}
Principled exception to the discrete-only stance: the closed Lie algebra $\mathfrak{so}(8) \supseteq G_2\cup\{\Delta\}$ generated from the Fano octonion left-multiplication table and the phase-lift generator $\Delta\in(\mathrm{e}_1,\mathrm{e}_7)$.  Although $\mathfrak{so}(8)$ is a continuous (28-dimensional real) Lie algebra, it is built here entirely from \emph{discrete} ingredients (the seven imaginary octonion units, a finite Fano table, and one $\pi/2$ rotation), and its closure to 28 dimensions is a finite linear-algebra certificate (see the \texttt{Hqiv.SO8Closure} / \texttt{Hqiv.SO8ClosureInterface} stack).  The Lie-group structure is therefore a \emph{de-facto continuum algebra obtained from a discrete construction}: continuous in parameter space, but with no continuum spacetime, no continuum measure, and no UV limit attached.  When this manuscript or its companions write $\exp(t\,\Delta)\in\mathrm{SO}(8)$ or treat $\mathfrak{so}(8)$ rotations as smooth, those are statements internal to a finite-dimensional Lie group whose generators are certified by discrete data; they are \emph{not} a continuum-spacetime commitment and do not require a continuum-QFT scaffolding.

% From papers/paper/*.tex: % Shared reader contract: patch QFT + discrete numerics.
% From papers/*.tex:     % Shared reader contract: patch QFT + discrete numerics.
% From papers/*.tex:     \input{include/patch_theory_messaging.tex}
% From papers/paper/*.tex: \input{../include/patch_theory_messaging.tex}

\paragraph{Patch QFT and discrete numerics (reader contract).}
HQIV (Horizon-Quantized Informational Vacuum) is formulated as a \emph{patch theory}: quantum fields, actions, and physical readouts are attached to discrete patches---shell indices $m\in\mathbb{N}$, finite spacetime index types (e.g.\ $\mathrm{Fin}\,4$), and horizon-limited charts---rather than to a hypothetical fundamental smooth spacetime obtained as a mesh-refinement or ultraviolet limit.
Accordingly, when this text refers to ``QFT,'' it means \emph{patch QFT}: patching rules, finite measurement ledgers, and variational identities on the discrete spine; it is not the claim that the theory is \emph{defined} by recovering ordinary continuum quantum field theory through such a limit.
The numbers that admit the sharpest statements---mass and coupling ladders, curvature ratios, shell thermodynamics, and rational witnesses in \texttt{HQIV\_LEAN}---are objects of \emph{discrete mathematics} on $\mathbb{N}$ (combinatorial growth, cumulative simplex ledgers) together with the ordered field $\mathbb{R}$ as used in the proof assistant for analysis, bounds, and phenomenological display.
Smooth-manifold or continuum field notation, when it appears, is a \emph{translation layer} for comparison with standard physics literature; it does not supersede the discrete definitions.

\paragraph{A continuum is not required, and in places would destroy these results.}
The discrete patch layer is \emph{not} a numerical approximation to a more fundamental continuum theory awaiting future construction; it is the ontology.  Where the literature treats ``the continuum limit'' as the goal, HQIV treats it as a \emph{calculation approximation} for comparison with standard formulas.  The continuum form is an optional convenience for comparison; the proved discrete statement is what carries the physics.  In several load-bearing places---the curvature imprint $\delta_E(m)$ at shell $m$, the harmonic ladder masses $m_\tau\to m_\mu\to m_e$, the lock-in vev $v=\sqrt{\eta_{\rm paper}\,\Omega_k(m_{\rm lockin};m_{\rm lockin})}$, the proton anchor at $m=\mathrm{referenceM}=4$, the baryon-asymmetry $\eta\propto 6^7\sqrt{3}/(m+1)$, and the rational coefficients $\alpha=3/5$, $\gamma=2/5$, $\alpha_{\rm GUT}=1/42$---taking a literal $m\to\infty$ or sub-Planck continuum limit would \emph{erase} the discrete index that is doing the predictive work.  Continuum forms of these objects are therefore not preferred refinements; they are degenerate, information-losing readouts of the same patch data.

\paragraph{An exception: $\mathfrak{so}(8)$ as a de-facto continuum algebra from a discrete construction.}
Principled exception to the discrete-only stance: the closed Lie algebra $\mathfrak{so}(8) \supseteq G_2\cup\{\Delta\}$ generated from the Fano octonion left-multiplication table and the phase-lift generator $\Delta\in(\mathrm{e}_1,\mathrm{e}_7)$.  Although $\mathfrak{so}(8)$ is a continuous (28-dimensional real) Lie algebra, it is built here entirely from \emph{discrete} ingredients (the seven imaginary octonion units, a finite Fano table, and one $\pi/2$ rotation), and its closure to 28 dimensions is a finite linear-algebra certificate (see the \texttt{Hqiv.SO8Closure} / \texttt{Hqiv.SO8ClosureInterface} stack).  The Lie-group structure is therefore a \emph{de-facto continuum algebra obtained from a discrete construction}: continuous in parameter space, but with no continuum spacetime, no continuum measure, and no UV limit attached.  When this manuscript or its companions write $\exp(t\,\Delta)\in\mathrm{SO}(8)$ or treat $\mathfrak{so}(8)$ rotations as smooth, those are statements internal to a finite-dimensional Lie group whose generators are certified by discrete data; they are \emph{not} a continuum-spacetime commitment and do not require a continuum-QFT scaffolding.

% From papers/paper/*.tex: % Shared reader contract: patch QFT + discrete numerics.
% From papers/*.tex:     \input{include/patch_theory_messaging.tex}
% From papers/paper/*.tex: \input{../include/patch_theory_messaging.tex}

\paragraph{Patch QFT and discrete numerics (reader contract).}
HQIV (Horizon-Quantized Informational Vacuum) is formulated as a \emph{patch theory}: quantum fields, actions, and physical readouts are attached to discrete patches---shell indices $m\in\mathbb{N}$, finite spacetime index types (e.g.\ $\mathrm{Fin}\,4$), and horizon-limited charts---rather than to a hypothetical fundamental smooth spacetime obtained as a mesh-refinement or ultraviolet limit.
Accordingly, when this text refers to ``QFT,'' it means \emph{patch QFT}: patching rules, finite measurement ledgers, and variational identities on the discrete spine; it is not the claim that the theory is \emph{defined} by recovering ordinary continuum quantum field theory through such a limit.
The numbers that admit the sharpest statements---mass and coupling ladders, curvature ratios, shell thermodynamics, and rational witnesses in \texttt{HQIV\_LEAN}---are objects of \emph{discrete mathematics} on $\mathbb{N}$ (combinatorial growth, cumulative simplex ledgers) together with the ordered field $\mathbb{R}$ as used in the proof assistant for analysis, bounds, and phenomenological display.
Smooth-manifold or continuum field notation, when it appears, is a \emph{translation layer} for comparison with standard physics literature; it does not supersede the discrete definitions.

\paragraph{A continuum is not required, and in places would destroy these results.}
The discrete patch layer is \emph{not} a numerical approximation to a more fundamental continuum theory awaiting future construction; it is the ontology.  Where the literature treats ``the continuum limit'' as the goal, HQIV treats it as a \emph{calculation approximation} for comparison with standard formulas.  The continuum form is an optional convenience for comparison; the proved discrete statement is what carries the physics.  In several load-bearing places---the curvature imprint $\delta_E(m)$ at shell $m$, the harmonic ladder masses $m_\tau\to m_\mu\to m_e$, the lock-in vev $v=\sqrt{\eta_{\rm paper}\,\Omega_k(m_{\rm lockin};m_{\rm lockin})}$, the proton anchor at $m=\mathrm{referenceM}=4$, the baryon-asymmetry $\eta\propto 6^7\sqrt{3}/(m+1)$, and the rational coefficients $\alpha=3/5$, $\gamma=2/5$, $\alpha_{\rm GUT}=1/42$---taking a literal $m\to\infty$ or sub-Planck continuum limit would \emph{erase} the discrete index that is doing the predictive work.  Continuum forms of these objects are therefore not preferred refinements; they are degenerate, information-losing readouts of the same patch data.

\paragraph{An exception: $\mathfrak{so}(8)$ as a de-facto continuum algebra from a discrete construction.}
Principled exception to the discrete-only stance: the closed Lie algebra $\mathfrak{so}(8) \supseteq G_2\cup\{\Delta\}$ generated from the Fano octonion left-multiplication table and the phase-lift generator $\Delta\in(\mathrm{e}_1,\mathrm{e}_7)$.  Although $\mathfrak{so}(8)$ is a continuous (28-dimensional real) Lie algebra, it is built here entirely from \emph{discrete} ingredients (the seven imaginary octonion units, a finite Fano table, and one $\pi/2$ rotation), and its closure to 28 dimensions is a finite linear-algebra certificate (see the \texttt{Hqiv.SO8Closure} / \texttt{Hqiv.SO8ClosureInterface} stack).  The Lie-group structure is therefore a \emph{de-facto continuum algebra obtained from a discrete construction}: continuous in parameter space, but with no continuum spacetime, no continuum measure, and no UV limit attached.  When this manuscript or its companions write $\exp(t\,\Delta)\in\mathrm{SO}(8)$ or treat $\mathfrak{so}(8)$ rotations as smooth, those are statements internal to a finite-dimensional Lie group whose generators are certified by discrete data; they are \emph{not} a continuum-spacetime commitment and do not require a continuum-QFT scaffolding.


\paragraph{Patch QFT and discrete numerics (reader contract).}
HQIV (Horizon-Quantized Informational Vacuum) is formulated as a \emph{patch theory}: quantum fields, actions, and physical readouts are attached to discrete patches---shell indices $m\in\mathbb{N}$, finite spacetime index types (e.g.\ $\mathrm{Fin}\,4$), and horizon-limited charts---rather than to a hypothetical fundamental smooth spacetime obtained as a mesh-refinement or ultraviolet limit.
Accordingly, when this text refers to ``QFT,'' it means \emph{patch QFT}: patching rules, finite measurement ledgers, and variational identities on the discrete spine; it is not the claim that the theory is \emph{defined} by recovering ordinary continuum quantum field theory through such a limit.
The numbers that admit the sharpest statements---mass and coupling ladders, curvature ratios, shell thermodynamics, and rational witnesses in \texttt{HQIV\_LEAN}---are objects of \emph{discrete mathematics} on $\mathbb{N}$ (combinatorial growth, cumulative simplex ledgers) together with the ordered field $\mathbb{R}$ as used in the proof assistant for analysis, bounds, and phenomenological display.
Smooth-manifold or continuum field notation, when it appears, is a \emph{translation layer} for comparison with standard physics literature; it does not supersede the discrete definitions.

\paragraph{A continuum is not required, and in places would destroy these results.}
The discrete patch layer is \emph{not} a numerical approximation to a more fundamental continuum theory awaiting future construction; it is the ontology.  Where the literature treats ``the continuum limit'' as the goal, HQIV treats it as a \emph{calculation approximation} for comparison with standard formulas.  The continuum form is an optional convenience for comparison; the proved discrete statement is what carries the physics.  In several load-bearing places---the curvature imprint $\delta_E(m)$ at shell $m$, the harmonic ladder masses $m_\tau\to m_\mu\to m_e$, the lock-in vev $v=\sqrt{\eta_{\rm paper}\,\Omega_k(m_{\rm lockin};m_{\rm lockin})}$, the proton anchor at $m=\mathrm{referenceM}=4$, the baryon-asymmetry $\eta\propto 6^7\sqrt{3}/(m+1)$, and the rational coefficients $\alpha=3/5$, $\gamma=2/5$, $\alpha_{\rm GUT}=1/42$---taking a literal $m\to\infty$ or sub-Planck continuum limit would \emph{erase} the discrete index that is doing the predictive work.  Continuum forms of these objects are therefore not preferred refinements; they are degenerate, information-losing readouts of the same patch data.

\paragraph{An exception: $\mathfrak{so}(8)$ as a de-facto continuum algebra from a discrete construction.}
Principled exception to the discrete-only stance: the closed Lie algebra $\mathfrak{so}(8) \supseteq G_2\cup\{\Delta\}$ generated from the Fano octonion left-multiplication table and the phase-lift generator $\Delta\in(\mathrm{e}_1,\mathrm{e}_7)$.  Although $\mathfrak{so}(8)$ is a continuous (28-dimensional real) Lie algebra, it is built here entirely from \emph{discrete} ingredients (the seven imaginary octonion units, a finite Fano table, and one $\pi/2$ rotation), and its closure to 28 dimensions is a finite linear-algebra certificate (see the \texttt{Hqiv.SO8Closure} / \texttt{Hqiv.SO8ClosureInterface} stack).  The Lie-group structure is therefore a \emph{de-facto continuum algebra obtained from a discrete construction}: continuous in parameter space, but with no continuum spacetime, no continuum measure, and no UV limit attached.  When this manuscript or its companions write $\exp(t\,\Delta)\in\mathrm{SO}(8)$ or treat $\mathfrak{so}(8)$ rotations as smooth, those are statements internal to a finite-dimensional Lie group whose generators are certified by discrete data; they are \emph{not} a continuum-spacetime commitment and do not require a continuum-QFT scaffolding.


\paragraph{Patch QFT and discrete numerics (reader contract).}
HQIV (Horizon-Quantized Informational Vacuum) is formulated as a \emph{patch theory}: quantum fields, actions, and physical readouts are attached to discrete patches---shell indices $m\in\mathbb{N}$, finite spacetime index types (e.g.\ $\mathrm{Fin}\,4$), and horizon-limited charts---rather than to a hypothetical fundamental smooth spacetime obtained as a mesh-refinement or ultraviolet limit.
Accordingly, when this text refers to ``QFT,'' it means \emph{patch QFT}: patching rules, finite measurement ledgers, and variational identities on the discrete spine; it is not the claim that the theory is \emph{defined} by recovering ordinary continuum quantum field theory through such a limit.
The numbers that admit the sharpest statements---mass and coupling ladders, curvature ratios, shell thermodynamics, and rational witnesses in \texttt{HQIV\_LEAN}---are objects of \emph{discrete mathematics} on $\mathbb{N}$ (combinatorial growth, cumulative simplex ledgers) together with the ordered field $\mathbb{R}$ as used in the proof assistant for analysis, bounds, and phenomenological display.
Smooth-manifold or continuum field notation, when it appears, is a \emph{translation layer} for comparison with standard physics literature; it does not supersede the discrete definitions.

\paragraph{A continuum is not required, and in places would destroy these results.}
The discrete patch layer is \emph{not} a numerical approximation to a more fundamental continuum theory awaiting future construction; it is the ontology.  Where the literature treats ``the continuum limit'' as the goal, HQIV treats it as a \emph{calculation approximation} for comparison with standard formulas.  The continuum form is an optional convenience for comparison; the proved discrete statement is what carries the physics.  In several load-bearing places---the curvature imprint $\delta_E(m)$ at shell $m$, the harmonic ladder masses $m_\tau\to m_\mu\to m_e$, the lock-in vev $v=\sqrt{\eta_{\rm paper}\,\Omega_k(m_{\rm lockin};m_{\rm lockin})}$, the proton anchor at $m=\mathrm{referenceM}=4$, the baryon-asymmetry $\eta\propto 6^7\sqrt{3}/(m+1)$, and the rational coefficients $\alpha=3/5$, $\gamma=2/5$, $\alpha_{\rm GUT}=1/42$---taking a literal $m\to\infty$ or sub-Planck continuum limit would \emph{erase} the discrete index that is doing the predictive work.  Continuum forms of these objects are therefore not preferred refinements; they are degenerate, information-losing readouts of the same patch data.

\paragraph{An exception: $\mathfrak{so}(8)$ as a de-facto continuum algebra from a discrete construction.}
Principled exception to the discrete-only stance: the closed Lie algebra $\mathfrak{so}(8) \supseteq G_2\cup\{\Delta\}$ generated from the Fano octonion left-multiplication table and the phase-lift generator $\Delta\in(\mathrm{e}_1,\mathrm{e}_7)$.  Although $\mathfrak{so}(8)$ is a continuous (28-dimensional real) Lie algebra, it is built here entirely from \emph{discrete} ingredients (the seven imaginary octonion units, a finite Fano table, and one $\pi/2$ rotation), and its closure to 28 dimensions is a finite linear-algebra certificate (see the \texttt{Hqiv.SO8Closure} / \texttt{Hqiv.SO8ClosureInterface} stack).  The Lie-group structure is therefore a \emph{de-facto continuum algebra obtained from a discrete construction}: continuous in parameter space, but with no continuum spacetime, no continuum measure, and no UV limit attached.  When this manuscript or its companions write $\exp(t\,\Delta)\in\mathrm{SO}(8)$ or treat $\mathfrak{so}(8)$ rotations as smooth, those are statements internal to a finite-dimensional Lie group whose generators are certified by discrete data; they are \emph{not} a continuum-spacetime commitment and do not require a continuum-QFT scaffolding.


\section{Introduction}

The HQIV framework is built on a discrete null-lattice whose combinatorial
growth, Fano incidence structure, and horizon readouts already produce
three-generation rigidity, shell-surface mass ladders, rational coupling
ratios ($\alpha=3/5$, $\gamma=2/5$), and an effective continuous $\xi$
chart without ever positing a fundamental smooth spacetime. At the same
time, the topological literature contains a striking result, established
by Nielsen: any gauge theory whose U(1) sector is both complete—realising
every principal U(1)-bundle over an arbitrary paracompact base—and
charge-quantised must be formulated, up to homotopy equivalence of the
base and isomorphism of bundles, on the universal complex Hopf fibration
$S^1\to S^\infty\to\mathbb{CP}^\infty$ and its finite approximations
$S^1\to S^{2n+1}\to\mathbb{CP}^n$ \cite{NielsenTUFT2026}.

Nielsen further demonstrates that the Standard-Model gauge factors arise
as natural reductions along this nested shell hierarchy: the circular
$S^1$ fibre yields the U(1) electromagnetic sector, the first non-trivial
shell $S^3$ furnishes SU(2), and the subsequent shell $S^5$ accommodates
SU(3). Gravity emerges from the Kähler geometry of the base together with
torsion induced by the fibre, producing a structure analogous to
Einstein–Cartan theory in which the Levi-Civita connection is recovered
precisely when the torsional degrees of freedom are set to zero. On each
shell the generalised Beltrami operator $B = d|_\xi$, acting on the
contact distribution, is shown to be elliptic and essentially
self-adjoint, with a discrete spectrum that remains stable under
torsion perturbations by the Kato–Rellich theorem \cite{NielsenTUFT2026}.

Fibre-winding decomposition yields independent topological sectors whose
Gaussian functional determinants, regularised via spectral zeta functions,
generate intrinsic mass scales without additional dimensionful input
beyond a single electroweak reference. Fermion mixing matrices are
likewise derived from the intersection-form overlaps of admissible cycles
in the cohomology of $\mathbb{CP}^4$, with CP-violating phases arising
directly from fibre holonomy. Dynamics, finally, are traced to the
fluctuation spectrum of the topological action on the outermost shell
$S^9$. These results furnish a unified geometric origin for both the
particle spectrum and the gauge-gravity structure.

This paper presents a rigorous examination of the interface between
Nielsen's construction and the HQIV framework. It identifies the precise
respects in which the nested Hopf shells, the contact Beltrami spectrum,
the fibre-winding sectors, and the fibre-induced holonomy of the former
illuminate structures that HQIV had previously derived from discrete
combinatorial and variational principles. The examination is supported
by machine-checked formalizations in the \texttt{HQIV\_LEAN} library,
which record both the alignments and the chart-specific distinctions
necessary for phenomenological accuracy.

The central methodological point is the shared commitment to finite
approximations. TUFT works with the finite Hopf shells $S^{2n+1}\to
\mathbb{CP}^n$ and treats the infinite case as a direct limit; HQIV works
with finite shells $m\in\mathbb{N}$ truncated at the lock-in horizon
$\mathrm{referenceM}=4$ and treats the continuous $\xi$ chart as an
effective readout generated by the discrete data. In both cases a literal
continuum is not required and in several places would erase the very
indices that carry the physics.

\subsection{Limitations and Scope}

The present study confines itself to those aspects of Nielsen's
framework that admit direct comparison with structures already present
in the HQIV formalism—namely the finite Hopf-shell hierarchy, the
contact Beltrami spectrum, fibre-winding multiplicity, and the geometric
origin of effective mass and coupling scales. No claim is advanced that
the full universality theorem, the single-Fermi-constant derivation of
the entire spectrum, or the detailed construction of the topological
action on $S^9$ have been reproduced or endorsed. The Lean development
likewise remains partial: while the low-winding spectral alignments and
the typed substrate for contact data are machine-verified, higher-order
structures such as intersection-form mixing matrices and the explicit
fluctuation spectrum on $S^9$ are still at the level of rigorously
formulated open targets.

A useful complementary perspective is supplied by the analysis of three-dimensional causal growth on the HQIV null lattice \cite{hqiv-3d-growth-paper}. That work isolates the quadratic shell-counting law \(A(m)=4(m+2)(m+1)\) as the combinatorial signature of \(3\)D causal growth and packages a machine-checked conditional forcing theorem: once a shared-manifold rapidity bridge and an explicit \(\mathfrak{so}(8)\) span witness are assumed as hypotheses, common rapidity, positivity and divergence of the curvature integral, and reference normalisation follow. The paper is explicit that a pure causal partial order does not, by itself, select the octonion multiplication table or the full internal gauge structure; the octonionic completion is conditional on additional geometric data. The same work develops a spanning-tree model of causal monogamy that derives the curvature-imprint pair \((\alpha,\gamma)=(3/5,2/5)\) and identifies the Fano plane as the natural \(3\)D projective spine for the seven imaginary directions. These results supply both positive combinatorial content for the target discrete \(3\)-complex used in the typed Hopf-shell mapping and a clear methodological limit: even within the HQIV framework, certain topological features require further structure beyond pure causal growth. This conditional character provides a natural counterpoint to the stronger necessity claims advanced in Nielsen's construction and helps delineate the precise scope of any combined programme.

\section{The finite-approximation discipline}

TUFT's construction begins with the observation that the classifying
space for U(1)-bundles with discrete charge is the infinite complex
projective space $\mathbb{CP}^\infty$, with the universal bundle the
infinite Hopf fibration $S^1\to S^\infty\to\mathbb{CP}^\infty$. All
finite approximations $S^1\to S^{2n+1}\to\mathbb{CP}^n$ are admissible
truncations. The Standard-Model factors arise as successive reductions:
U(1) from the circular fiber, SU(2) from the $S^3$ shell, SU(3) from the
$S^5$ shell. Gravity emerges from the Kähler geometry of the base
together with torsion induced by the fiber.

HQIV never begins from a bundle classification. It begins from a discrete
null-lattice whose vertices are tagged by stars-and-bars at each shell
$m$, whose incidence is governed by the Fano plane, and whose horizon
areas $S(m)=(m+1)(m+2)$ supply the surfaces that appear in every mass
and coupling readout. Three generations appear because the integrable
fiber-winding sectors that survive in the discrete complex are exactly
three; the lock-in at $m=4$ is a combinatorial closure condition, not a
fitted vev.

The two pictures meet at the level of finite truncations. The TUFT
statement that only $n=1,2,3$ yield integrable torus sectors (unknot,
Hopf link, trefoil) before a hyperbolic transition at $n=4$ supplies a
topological reason for the HQIV fact that $\mathrm{ResonanceGeneration}$
is exactly $\mathrm{Fin}\,3$. The finite Hopf shells supply a reason why
a discrete shell ladder with a distinguished lock-in horizon should
exist at all. Conversely, the HQIV patch ontology supplies the rule that
neither $\mathbb{CP}^\infty$ nor the continuous $\xi$ chart may be
treated as fundamental; they are limiting readouts of finite data.

\section{Contributions of Nielsen's Topological Framework to HQIV}

The value of the TUFT perspective is not that it supplies new numerical
fits, but that it supplies a topological \emph{necessity} argument and a
spectral-geometric origin story for structures that HQIV had previously
derived from discrete combinatorial and variational axioms. The most
substantial contributions fall under four headings.

\subsection{Classifying-space necessity and indecomposability}

TUFT proves that any gauge theory whose U(1) sector is both complete
(realises every principal U(1)-bundle over paracompact bases) and
charge-quantised (admits only discrete charges) must, up to homotopy
equivalence and bundle isomorphism, be formulated on the universal
complex Hopf fibration and its finite approximations. The resulting
total structure group $G_{\mathrm{total}} = (\mathrm{SU}(3)\times
\mathrm{SU}(2)\times\mathrm{U}(1)\times\mathrm{SO}(4))/\Gamma$ is
intrinsically non-factorable because the first Chern class $c_1$
generates the cohomology ring $H^*(\mathbb{CP}^\infty;\mathbb{Z})\simeq
\mathbb{Z}[c_1]$.

In the HQIV setting this supplies a principled reason why the Fano
plane, the octonion carrier, and the three-generation count appear
together rather than as independent modelling choices. The discrete
null-lattice with its seven-line incidence structure can now be read as
a combinatorial model of the admissible reductions along the first few
Hopf shells. The non-factorability claim gives additional weight to the
already-proved algebraic closure results that generate $\mathfrak{so}(8)$
from $G_2\cup\{\Delta\}$ inside the same eight-dimensional carrier.

\subsection{Finite approximations and the ``no infinities'' discipline}

Both frameworks are committed to finite truncations. TUFT works with
$S^{2n+1}\to\mathbb{CP}^n$ and treats $\mathbb{CP}^\infty$ as a direct
limit; HQIV works with shells $m\in\mathbb{N}$ truncated at the lock-in
horizon $\mathrm{referenceM}=4$ and treats the continuous $\xi$ chart as
an effective readout. The TUFT statement that only windings $n=1,2,3$
yield integrable torus sectors before a hyperbolic transition at $n=4$
furnishes an independent topological reason for the HQIV fact that
$\mathrm{ResonanceGeneration}$ is exactly $\mathrm{Fin}\,3$. Conversely,
the HQIV insistence that continuum limits are information-destroying
readouts (rather than ontological improvements) supplies the
methodological safeguard that prevents the TUFT infinite fibration from
being misread as a commitment to a fundamental smooth spacetime.

\subsection{Contact spectrum, holonomy phases, and mixing}

TUFT's generalised Beltrami operator on the contact distribution of each
Hopf shell yields a discrete spectrum stable under torsion perturbations
(Kato--Rellich). The Lean formalisation in
\hyperref[lean:hopf-shell-complex]{HopfShellComplex} together with the
earlier bridges records both the shared eigenvalue law
$\lambda_\ell=\ell(\ell+2)$ on $S^3$ and the explicit normalisation
distinction between TUFT's fundamental coexact mode and the scalar
Peter--Weyl value. The fibre-winding decomposition supplies independent
topological sectors whose zeta-regularised Gaussian determinants are
intended to generate intrinsic mass scales; the existing HQIV
informational-energy and sector-Gaussian-leading-weight constructions
are now seen as discrete, patch-level realisations of precisely this
mechanism.

Fermion mixing (CKM/PMNS) is claimed to arise from intersection-form
overlaps of admissible cycles in $H^*(\mathbb{CP}^4)$, with CP violation
induced by fibre holonomy phases. The rh-fourier-lift curvature channel
$K(n,\alpha)$ and its associated $\mathrm{PhaseMap}$ (already wired to
the new Hopf-shell holonomy carrier) provide a concrete discrete
realisation of those phases. The existing Fano holonomy rows and
imprint phase machinery therefore acquire a possible topological
interpretation rather than remaining purely phenomenological charts.

\subsection{Single-scale derivation versus multi-witness architecture}

TUFT asserts that a single empirical input—the Fermi constant together
with its associated electroweak vacuum expectation value—suffices to
determine the fine-structure constant and all shell-specific mass
scales, spectral coefficients, and couplings via the spectral geometry
of the fibration. HQIV, by contrast, maintains an explicitly
multi-witness architecture (proton lock-in, CODATA $\alpha$, present-CMB)
with operational discipline about which scale is active in any given
solve. Nielsen's framework therefore presents a research programme that
invites comparison with HQIV's current multi-witness architecture. The
machine-checked chart distinctions recorded in the Lean library (the
$4/3$ alignment at the lock-in neighbour under the $m = n + 1$ chart, the
location of the $3/2$ ratio within Fano holonomy rows rather than
geometric resonance steps, and the separation of lepton resonance ratios
from the low-winding Beltrami ratios) provide the necessary precision for
any such comparison.

The new typed layer \texttt{Hqiv.Topology.HopfShellComplex} (with its
functor to the discrete 3-complex substrate and its PhaseMap carrier
hook) makes several of the above alignments machine-checkable rather
than narrative. It does not yet reproduce the full contact geometry,
Ray--Singer torsion, or the universality theorem; those remain
explicitly open targets.

\subsection{Re-interpretation of TUFT's dark sector}

TUFT lists ``dark sector effects from bundle torsion and holonomy'' among
the consequences of its construction. HQIV has no separate dark sector
by construction: all energy-density contributions are internal to the
informational vacuum of the discrete null-lattice plus its curvature and
phase structure.

The precise correspondence is as follows. TUFT derives mass scales and
apparent extra contributions from the zeta-regularised Gaussian
determinants of its fibre-winding sectors together with torsion and
holonomy corrections on the Hopf fibration. In HQIV the dimensioned
quantities observed at the present epoch (``now'') are given by the
informational energy
\[
E_{\mathrm{tot}}(m_{\mathrm{rest}},\xi_{\mathrm{now}})
= m_{\mathrm{rest}} + \frac{1}{\Theta_{\mathrm{local}}(\xi_{\mathrm{now}})},
\qquad
\Theta_{\mathrm{local}}(\xi) = \frac{T_{\mathrm{Pl}}}{\xi},
\]
together with additive corrections from the curvature channel
$K(n,\alpha)$ (whose $\alpha\log x$ term supplies the imprintable
geometric contribution) and the rapidity/phase lifts already wired to
the Hopf-shell holonomy carrier.

Thus TUFT's ``derived constants'' (mass scales, effective couplings,
spectral coefficients obtained from a single Fermi-scale input via Hopf
spectral geometry) are exactly the effective values of HQIV's
dimensioned now quantities evaluated at the current horizon coordinate
$\xi_{\mathrm{now}}$ (or the corresponding lock-in or outer-horizon
shell in the discrete ladder), including the informational localization
and curvature-imprint terms. What appears in TUFT as a dark sector is,
under the HQIV reading, the internal energetic contribution of the
informational vacuum itself once it is dimensionally evaluated at now.
No additional fields or sectors are required.

A concrete experimental discriminant arises from the contrasting origins
assigned to apparent new effects. Nielsen's framework predicts a
torsion-induced phase wobble generated by fibre holonomy, detectable in
precision interferometry or spin-precession measurements. The HQIV
account instead attributes analogous signatures to the curvature-channel
imprint: the $\alpha\log x$ term in the discrete density $\rho(x)$
accumulates into $K(n,\alpha)$ and modulates effective couplings at the
present epoch. Distinguishing global fibre-holonomy wobble from local
curvature-imprint modulation therefore provides a sharp test accessible
to existing and near-future precision experiments. A quantitative
estimate of the expected wobble amplitude (for example in
current-generation atom interferometers or spin-precession measurements)
would render the discriminant still more incisive; such an estimate is
not yet available within the present formalisation and is reserved for
subsequent work.

This re-interpretation sharpens the single-scale versus multi-witness
comparison: the additional contributions that must be derived
geometrically in Nielsen's framework appear in HQIV as the explicit
$1/\Theta$ localisation and curvature-channel terms already present in
the definition of observable quantities at the present epoch. The Lean
formalisation of the informational-energy readouts and the curvature
channel $K(n,\alpha)$ therefore provides a direct computational bridge
for quantitative comparison.

\section{Proved Lean alignments (the honest bridge)}

All alignments below are stated as theorems that have been formally
verified in Lean without the use of `sorry`. They are deliberately
chart-specific; the library contains explicit mismatch theorems wherever
a naïve identification would be false.

\subsection{Beltrami spectrum and normalization distinction}

The scalar Laplace--Beltrami eigenvalues on the unit $S^3$ are
$\lambda_\ell=\ell(\ell+2)$. TUFT's coexact Beltrami operator on the
contact distribution yields the same law. The Lean library records both
normalizations and proves they differ at the fundamental level:

\begin{itemize}
\item $\mathrm{beltramiPeterWeylEigenvalueS3}\,\ell =
  \mathrm{laplaceBeltramiEigenvalueS3}\,\ell$ (shared law).
\item The TUFT fundamental coexact mode uses $\lambda_1=2$; the Peter--Weyl
  level $\ell=1$ gives 3. The theorem
  $\mathrm{tuftFundamentalBeltrami\_ne\_eq\_peterWeyl\_one}$ records the
  mismatch explicitly.
\end{itemize}

Multiplicity on the $S^3$ shell is $(n+1)^2$ for winding $n$, matching
TUFT sector multiplicity.

\subsection{Three generations from integrable fiber windings}

TUFT states that only the first three positive fiber windings yield
integrable torus sectors. The Lean definition
\begin{quote}
$\mathrm{HopfFiberWinding}\,n \;\;:\;\; \mathrm{Prop}$
\end{quote}
is true precisely for $n=1,2,3$ and false for $n=0$ and $n\ge4$. The
theorems $\mathrm{hopfIntegrableGenerationCount\_eq\_three}$ and
$\mathrm{tuftMinimalBeltrami\_matches\_fin3}$ establish that this
counting is identical to the HQIV $\mathrm{ResonanceGeneration}=\mathrm{Fin}\,3$
and that the Beltrami eigenvalues are strictly ordered on those three
slots.

\subsection{Lock-in neighbour match and chart distinctions}

Under the bookkeeping chart $m=n+1$ (TUFT winding to HQIV discrete shell),
the minimal Beltrami ratio for windings $3\to2$ is exactly $4/3$. On the
HQIV side the geometric resonance step at the lock-in neighbour shells is
\begin{quote}
$\mathrm{geometricResonanceStep}\,4\,3 = 4/3$.
\end{quote}
The equality
$\mathrm{tuftBeltrami\_tau\_mu\_eq\_geometricResonance\_lockin\_neighbor}$
is a proved theorem.

The same machinery records the mismatch at the previous step: the Beltrami
ratio $2\to1$ is $3/2$, but the geometric step on shells $3\to2$ is
$35/24$. The $3/2$ ratio does appear, but in the Fano holonomy rows
($\mathrm{holonomyRowRhs}$ on the heavy and middle vertices) rather than
in the geometric resonance steps. These distinctions are not bugs; they
are the content of the honest bridge.

\subsection{Informational energy and spectral weight scaffold}

TUFT obtains intrinsic scales from zeta-regularized Gaussian determinants
over the fiber-winding sectors. HQIV already possesses an informational
energy
\[
E_{\mathrm{tot}}(m_{\mathrm{rest}},\xi) = m_{\mathrm{rest}} +
\frac{1}{\Theta(\xi)}
\]
and a Beltrami-level correction term
$\mathrm{beltramiSpectralWeightS3}\,\ell =
(\mathrm{beltramiPeterWeylEigenvalueS3}\,\ell + 1)^{-1}$.
The scaffold $\mathrm{informationalEnergyAtXiWithBeltrami}$ packages the
sum. The deeper identification of this term with a per-sector zeta
determinant remains open (per-sector zeta target).

\section{Discrete topology pegs already present}

The library contains a discrete 3-complex layer
(\texttt{Hqiv/Topology/DiscreteNullLatticeComplex.lean}) with
$\mathrm{NullShellVertex}$, $\mathrm{Discrete3Complex}$, Euler
characteristic, combinatorial sphericity ($\chi=0$ for $S^3$-like
references), and an $\mathrm{S3NullReference}$ template whose
edges and triangles are currently empty. This is the natural host for
Hopf-shell contact data or fiber-winding sectors. The new
\hyperref[lean:hopf-shell-complex]{Hopf-shell complex module} supplies the
typed `HopfShell` + `ContactBeltrami` layer that maps directly onto this
substrate (Hopf-shell complex + ContactBeltrami core).

Parallel modules (\texttt{ParallelPoincareScaffold}, \texttt{DiscretePhaseEvolution},
\texttt{SignedShellBudget}) already carry discrete Poincaré and phase-evolution
language compatible with torsion and holonomy without continuum commitments.

\section{Synergy with the rh-fourier-lift curvature channel}

On the current development branch the curvature density
$\rho(x,\alpha)=(1+\alpha\log x)/x$ and its cumulative channel $K(n,\alpha)$
are proved to dominate the harmonic series, to be strictly monotone for
$\alpha>0$, and to admit a phase lift $\mathrm{PhaseMap}$ (canonical
example: rigid rotation by $\log\omega$). This discrete phase-lift
mechanism is the natural carrier for TUFT fiber holonomy phases. The
same channel supplies a concrete, finite-n model for the torsion that TUFT
derives from the fiber in the Einstein--Cartan-like sector.

\section{Exact trapped-Casimir factorization of the binding cell}

A major recent formal advance is the explicit algebraic factorization of
the existing HQIV binding law as a lattice-counted trapped zero-point
(Casimir) energy term plus one normalized selection factor drawn from the
SO(8) carrier:

\[
E_{\mathrm{bind}} = \text{lattice count} \times
(\text{per-mode zero-point}) \times
\text{normalized SO(8) trace selection}.
\]

In Lean this is recorded by the definitions
`normalizedSO8TraceSelection` and `trappedCasimirCouplingCell`, with
proved identities
`trappedCasimirCouplingCell_eq_alphaEffAtShell` and
`bindingCouplingAtShell_eq_lattice_trappedCasimirCell`
(and the equivalent form using `available_modes`). The factorization is
exact; no continuum limit or external gluon field is required for the
algebra.

The phase-lift and three-shell layer now supplies a named closure predicate
\leanid{T11T12ClosesTrappedCasimirSelection} (a Hopf shell's torsion coefficient
supplies exactly the required normalized selection at a given readout
shell) together with a conditional theorem that turns the ontology into a
precise proof obligation. The concrete three-shell non-factorability witness is
wired directly into this coefficient.

This work does not yet claim that ``gluons do not exist.'' It isolates the
residual as a sharply defined selection factor on the closed carrier and
makes the statement ``the phase-lift and inner-shell torsion data on the integrable Hopf
shells reproduces that selection'' a theorem target rather than a
narrative hope. The per-shell `curvatureImprintAlpha` hook (with the
associated `trappingSelectionFromHeavyHopfShellWithAlpha` and three-shell
variants) supplies the natural geometric dial for exploring whether
different contact geometries on the successive Hopf shells can furnish
different effective imprints that move both stabilization horizons and
trapping factors.

\section{Open Phase-2 targets}

Substantial machine-checked progress has been made on the Phase-2 topology
program (typed Hopf-shell complexes with direct mapping into the
null-lattice 3-complex, ContactBeltrami spectrum and formal stability
scaffolding, per-sector zeta-determinant leading-term witnesses, the
three-row phase-contribution assembler with full three-row coverage including the
light slot together with admissible-cycle and orientation fields, canonical
matrix torsion with skew-adjoint action and Parallel-Poincaré
bridge, and the concrete three-shell non-factorability witness with closure condition).
The exact trapped-Casimir factorization of the binding cell (see above)
and the active per-shell curvature-imprint program
(`curvatureImprintAlpha` hook + explicit `trappingSelection*WithAlpha`
objects) constitute the current research frontier.

The early mass-spectrum formalization targets (lepton chart through lock-in
propagation) remain partly open, though the lepton-chart layer now possesses
strong interlocks with the phase assembler and inner three-shell witness; the MeV chart layer has been made fully
parametric. The deeper outer neutral-shell and $S^9$ fluctuation program and several
Phase-3 topics (real admissible-cycle overlap forms, genuine CP-phase
threading, full contact operator on forms) are still at the scaffold or
early-theorem stage. They are recorded here only as finite-patch development
directions; none of them import TUFT's universality claim or its single-input
derivation of the full spectrum.

\section{Conclusion}

TUFT supplies a topological derivation for the existence and structure
of the discrete shells, the three-generation count, the contact spectrum,
and the fiber holonomy that appear in HQIV as geometric or combinatorial
facts. HQIV supplies the patch-ontology discipline and the finite-truncation
rule that keep both pictures from treating their continuous limiting
objects as fundamental. The Lean library already certifies the precise
points of alignment and the precise points of chart mismatch. The deeper
program---populating the discrete complex with Hopf data, lifting the
contact Beltrami operator, and exhibiting the curvature channel as a
discrete torsion source---is now a well-scoped, bidirectional research
program rather than a narrative hope.

\appendix
\section{Lean module index}\label{app:lean-index}

The following reader-friendly names appear as hyperlinked anchors in the
body. Each entry records the Lean module and the principal identifier(s).

\begin{description}

\item[\phantomsection\label{lean:beltrami-peter-weyl}Beltrami--Peter--Weyl eigenvalue law.]
  Shared spectral law on unit $S^3$. Lean:
  \leanid{beltramiPeterWeylEigenvalueS3},
  \leanid{beltramiPeterWeylEigenvalueS3\_eq\_laplace}
  in \leanfile{Hqiv/Physics/HopfShellBeltramiMassBridge.lean}.

\item[\phantomsection\label{lean:tuft-fundamental-ne}TUFT fundamental vs.\ Peter--Weyl normalization.]
  Explicit mismatch at the fundamental level. Lean:
  \leanid{tuftFundamentalBeltrami\_ne\_eq\_peterWeyl\_one}
  in \leanfile{Hqiv/Physics/HopfShellBeltramiMassBridge.lean}.

\item[\phantomsection\label{lean:hopf-fiber-winding}Hopf fiber winding (integrable sectors).]
  Exactly three positive windings yield integrable torus sectors.
  Lean: \leanid{HopfFiberWinding},
  \leanid{hopfIntegrableGenerationCount\_eq\_three},
  \leanid{tuftMinimalBeltrami\_matches\_fin3}
  in \leanfile{Hqiv/Physics/HopfShellBeltramiMassBridge.lean}.

\item[\phantomsection\label{lean:tuft-minimal-beltrami}Minimal Beltrami eigenvalue per winding.]
  $\lambda_{\mathrm{min}}(n)=n+1$. Lean:
  \leanid{tuftMinimalBeltramiEigenvalue}
  and the one/two/three evaluation theorems
  in \leanfile{Hqiv/Physics/HopfShellBeltramiMassBridge.lean}.

\item[\phantomsection\label{lean:tuft-beltrami-ratios}TUFT Beltrami resonance ratios.]
  $3\to2$ gives $4/3$; $2\to1$ gives $3/2$. Lean:
  \leanid{tuftBeltramiResonanceRatio\_tau\_mu},
  \leanid{tuftBeltramiResonanceRatio\_mu\_e}
  in \leanfile{Hqiv/Physics/HopfShellBeltramiMassBridge.lean}.

\item[\phantomsection\label{lean:lockin-4-3-match}Lock-in neighbour $4/3$ alignment.]
  Under the $m=n+1$ chart the Beltrami ratio equals the geometric
  resonance step at the lock-in neighbour shells. Lean:
  \leanid{tuftBeltrami\_tau\_mu\_eq\_geometricResonance\_lockin\_neighbor},
  \leanid{geometricResonanceStep\_four\_three\_eq\_four\_thirds}
  in \leanfile{Hqiv/Physics/HalfStepBeltramiShellBridge.lean}.

\item[\phantomsection\label{lean:chart-mismatch-3-2}Explicit chart mismatch at $2\to1$.]
  Beltrami $3/2$ is not the geometric step $35/24$ on shells $3\to2$.
  Lean: \leanid{tuftBeltrami\_mu\_e\_ne\_geometricResonanceStep\_three\_two}
  in \leanfile{Hqiv/Physics/HalfStepBeltramiShellBridge.lean}.

\item[\phantomsection\label{lean:holonomy-3-2-placement}3/2 ratio in Fano holonomy rows.]
  The Beltrami $2\to1$ ratio equals the holonomy-row budget ratio on
  the heavy and middle Fano vertices. Lean:
  \leanid{tuftBeltrami\_mu\_e\_eq\_holonomyRowRhs\_vertices},
  \leanid{tuftBeltrami\_mu\_e\_eq\_fanoHolonomyWeight\_ratio}
  in \leanfile{Hqiv/Physics/HalfStepBeltramiShellBridge.lean}.

\item[\phantomsection\label{lean:sector-gaussian-weight}Sector Gaussian leading weight.]
  $\mathrm{detunedShellSurface} = \mathrm{shellSurface} /
  \mathrm{omaxwellFanoDetuning1Jet}$. Lean:
  \leanid{sectorGaussianLeadingWeight\_eq\_detunedShellSurface}
  in \leanfile{Hqiv/Physics/FanoSectorSpectralMassEmergence.lean}.

\item[\phantomsection\label{lean:laplace-s3-s4}Laplace--Beltrami on $S^3$ and $S^4$.]
  Eigenvalue laws and dimension formulas. Lean:
  \leanfile{Hqiv/Geometry/QuaternionMaxwellS3OMaxwellS4Spectral.lean}.

\item[\phantomsection\label{lean:discrete-3-complex}Discrete 3-complex layer.]
  $\mathrm{NullShellVertex}$, $\mathrm{Discrete3Complex}$,
  Euler characteristic, combinatorial sphericity, $\mathrm{S3NullReference}$
  template. Lean:
  \leanfile{Hqiv/Topology/DiscreteNullLatticeComplex.lean}.

\item[\phantomsection\label{lean:curvature-channel}Curvature channel and phase map (rh-fourier-lift).]
  $\rho(x,\alpha)$, $K(n,\alpha)$ with proved domination and monotonicity;
  $\mathrm{PhaseMap}$ abstraction with canonical log-$\omega$ lift.
  Lean: \leanfile{RhFourierLift/Setup.lean}.

\item[\phantomsection\label{lean:hopf-shell-complex}Hopf-shell complex + ContactBeltrami (topology core).]
  Typed `HopfShell` with integrable-winding certificate, functor
  `toDiscrete3Complex_integrable` landing in the null-lattice
  3-complex substrate, `ContactBeltrami` spectrum and multiplicity
  agreement with existing Beltrami data, and `HolonomyPhaseCarrier`
  hook for rh-fourier-lift `PhaseMap`. Lean:
  \leanfile{Hqiv/Topology/HopfShellComplex.lean}. Now part of the
  official `HQIVParallelPoincare` and `HQIVPhysics` build globs;
  imported by `ShellOpeningEvolution` (shell-opening wiring) and the Beltrami
  bridge (compatibility transport).

\item[\phantomsection\label{lean:informational-energy-beltrami}Informational energy with Beltrami correction.]
  Scaffold adding the inverse Beltrami weight to the localization term.
  Lean: \leanid{informationalEnergyAtXiWithBeltrami},
  \leanid{beltramiSpectralWeightS3}
  in \leanfile{Hqiv/Physics/HopfShellBeltramiMassBridge.lean}.

\item[\phantomsection\label{lean:trapping-selection}Trapping selection with per-shell imprints.]
  Explicit objects for exploring different effective curvature imprints
  $\alpha_n$ on the integrable Hopf shells.
  Lean: \leanid{trappingSelectionFromHeavyHopfShell},
  \leanid{trappingSelectionFromHeavyHopfShellWithAlpha},
  \leanid{trappingSelectionFromThreeHopfShellsWithAlphas}
  in \leanfile{Hqiv/Physics/HopfShellBeltramiMassBridge.lean}.

\item[\phantomsection\label{lean:trapped-casimir-factorization}Exact trapped-Casimir factorization of binding.]
  Algebraic factorization of the binding cell as lattice-counted
  zero-point energy plus normalized SO(8) selection.
  Lean: \leanid{normalizedSO8TraceSelection},
  \leanid{trappedCasimirCouplingCell},
  \leanid{trappedCasimirCouplingCell_eq_alphaEffAtShell},
  \leanid{bindingCouplingAtShell_eq_lattice_trappedCasimirCell},
  \leanid{T11T12ClosesTrappedCasimirSelection},
  \leanid{bindingCouplingAtShell_eq_T11T12_trappedCasimir}
  in \leanfile{Hqiv/Physics/HopfShellBeltramiMassBridge.lean}.

\end{description}

\bibliographystyle{plain}
\bibliography{../references}

\end{document}